% TOP_PROBLEM: {"id": 30007041, "title": "A geometric criterion for Kuga–Satake algebraicity", "category_id": 5, "category": {"id": 5, "name": "algebraic_geometry", "display_name": "Algebraic Geometry", "description": "Geometric objects defined by polynomial equations.", "slug": "algebraic-geometry", "order_index": 5, "created_at": "2026-07-31T15:26:25.671Z"}, "status": "open"}
% FIRST_POSTED: 2026-10-06
% ATTEMPT_STATUS: solved
% OPENAI_CLAIM: {"schema_version":1,"source_repo":"https://github.com/openai/math","source_commit":"adc7f1241b42e322a6451854ab7e4b4c146bf78a","release_date":"2026-10-06","match":"full","resolution":"proved","independently_reviewed":false,"second_pass":"agreed","families":[{"family":"032","title":"Hodge and Kuga–Satake results for all projective K3 surfaces","subject":"Algebraic and complex geometry","manuscripts":[{"title":"A Conditional Reduction for Algebraic Kuga–Satake Correspondences","pdf_path":"preprints/A-Conditional-Reduction-for-Algebraic-Kuga-Satake-Correspondences-September-10-2026/paper.pdf","readme_path":"preprints/A-Conditional-Reduction-for-Algebraic-Kuga-Satake-Correspondences-September-10-2026/README.md","date":"2026-09-10","theorem_ref":"main:generic (sections/introduction.tex:36)","statement_sources":[{"path":"preprints/A-Conditional-Reduction-for-Algebraic-Kuga-Satake-Correspondences-September-10-2026/build/sections/introduction.tex","line":36,"label":"main:generic"}]}],"lean":null,"reasoning_trace":null}]}
\documentclass[11pt]{article}
\usepackage[margin=25mm]{geometry}
\usepackage{fontspec}
\setmainfont{Latin Modern Roman}
\usepackage{amsmath,amssymb,mathtools}
\usepackage{unicode-math}
\setmathfont{Latin Modern Math}
\usepackage{enumitem}
\usepackage{xurl,hyperref,bookmark}
\hypersetup{colorlinks=true,urlcolor=blue}
\setcounter{secnumdepth}{0}
\setlength{\parindent}{0pt}
\setlength{\parskip}{5pt}
\setlength{\emergencystretch}{3em}
\begin{document}
\section*{A geometric criterion for Kuga–Satake algebraicity}

\subsection{OpenAI's claim}
A nonzero algebraic map from generic primitive K3 cohomology to an abelian target is claimed to recover the exact Kuga--Satake correspondence; existence of that starting map remains an assumption. The cited manuscript argues for that result; its relation to this problem is specified below.

\subsection{Mathematical statement of the claimed result}
\paragraph{Mathematical claim 1.} Let $(S,h)$ be a primitively amply polarized projective complex K3 surface, put $V=h^\perp\subset H^2(S,\mathbb Q)$ with cup-product form $q$, and assume $\operatorname{Hg}(V)=\operatorname{SO}(V,q)$. This hypothesis implies $T(S)=\operatorname{NS}(S)_{\mathbb Q}^{\perp}=V$. For any stipulated rational quadratic polarization and full standard even-Clifford Kuga--Satake construction, including parity, weight-one polarization and isogenous realization, write $A_S$ for its abelian variety and $\kappa_S:V\to H^2(A_S^2,\mathbb Q)$ for the prescribed map. OpenAI asserts equivalence of these three existence statements: (i) some abelian variety $B$ admits a rational codimension-two cycle on $S\times B$ inducing a nonzero map from $V$ to $H^2(B,\mathbb Q)$; (ii) a smooth connected projective surface $Y$ admits a dominant generically finite map $f:Y\to S$ with
\[
f_*\operatorname{im}\bigl(\bigwedge^2 H^0(Y,\Omega_Y^1)\longrightarrow H^0(Y,\Omega_Y^2)\bigr)\ne0;
\]
(iii) some $\Gamma\in\operatorname{CH}^2(S\times A_S^2)_{\mathbb Q}$ satisfies $\Gamma_*|_V=\kappa_S$. Thus algebraicity is conditional on either equivalent geometric input; this theorem does not assert that every K3 surface has that input.

\subsection{Outline of the argument}
Manuscript 1 (family 032, 2026-09-10): Acting independently on each weight-one factor reconstructs the Clifford multiplicities, and rational descent treats nonsplit algebras. Hilbert parameter spaces carry the specified flat cycle class to special periods, while signed square-root comparisons keep the polarization normalization fixed.

\subsection{Scope relative to this problem}
Manuscript 1 (family 032, 2026-09-10): The problem selects the mathematical statement summarized here, including its stated hypotheses and parameter range. The cited theorem gives a claimed proof of that assertion in full.

\subsection{Formal verification}
Family 032 has no Lean scope page in this snapshot; no formalization of these manuscript results is listed.

This site has not executed the Lean files or checked their dependencies. These descriptions report the scope of the linked declarations, rather than a new verification of the proofs.

\subsection{Status caveat}
OpenAI attributes these results to an internal model. Its release notice cautions that results lacking full formalization can contain mistakes. This site has not independently assessed the proofs; a solved label records the scope of the reported claim.

\subsection{Sources}
\begin{itemize}
\item OpenAI, \emph{A Conditional Reduction for Algebraic Kuga–Satake Correspondences}, 2026-09-10; citation key \texttt{OAI:A-Conditional-Reduction-for-Algebraic-Kuga-Satake-Correspondences-September-10-2026}. \href{https://github.com/openai/math/blob/adc7f1241b42e322a6451854ab7e4b4c146bf78a/preprints/A-Conditional-Reduction-for-Algebraic-Kuga-Satake-Correspondences-September-10-2026/paper.pdf}{PDF (GitHub)}; \href{https://github.com/openai/math/raw/adc7f1241b42e322a6451854ab7e4b4c146bf78a/preprints/A-Conditional-Reduction-for-Algebraic-Kuga-Satake-Correspondences-September-10-2026/paper.pdf}{PDF (direct)}; \href{https://github.com/openai/math/blob/adc7f1241b42e322a6451854ab7e4b4c146bf78a/preprints/A-Conditional-Reduction-for-Algebraic-Kuga-Satake-Correspondences-September-10-2026/README.md}{Manuscript page and citation}; \href{https://github.com/openai/math/blob/main/preprints/A-Conditional-Reduction-for-Algebraic-Kuga-Satake-Correspondences-September-10-2026/paper.pdf}{Latest version}; Statement source: \href{https://github.com/openai/math/blob/adc7f1241b42e322a6451854ab7e4b4c146bf78a/preprints/A-Conditional-Reduction-for-Algebraic-Kuga-Satake-Correspondences-September-10-2026/build/sections/introduction.tex\#L36}{main:generic}.
\item \href{https://github.com/openai/math/blob/adc7f1241b42e322a6451854ab7e4b4c146bf78a/CONTENTS.md}{OpenAI catalogue entry}.
\end{itemize}
\end{document}
